\documentclass[11pt]{article}

\usepackage[a4paper,margin=1in]{geometry}
\usepackage{amsmath,amssymb,amsthm,mathtools}
\usepackage{enumitem}
\usepackage{hyperref}

\newtheorem{theorem}{Theorem}
\newtheorem{proposition}{Proposition}
\newtheorem{definition}{Definition}
\newtheorem{example}{Example}
\newtheorem{remark}{Remark}

\DeclareMathOperator{\Law}{Law}
\DeclareMathOperator{\supp}{supp}
\DeclareMathOperator{\id}{id}

\newcommand{\R}{\mathbb R}
\newcommand{\E}{\mathbb E}
\newcommand{\Pcal}{\mathcal P}
\newcommand{\Borel}{\mathcal B}
\newcommand{\one}{\mathbf 1}
\newcommand{\dd}{\,\mathrm d}
\newcommand{\PiC}{\Pi}

\title{Stochastic Orders, Strassen's Theorem, and Optimal Transport}
\author{}
\date{}

\begin{document}

\maketitle

\section{Overview}

A stochastic order is an order relation between probability measures. Instead of
comparing two points \(x\) and \(y\), one compares two laws \(\mu\) and \(\nu\)
by asking whether \(\nu\) is larger, more variable, or more spread out than
\(\mu\), according to some chosen class of test functions.

Strassen's theorem is the bridge between two viewpoints:
\begin{itemize}[leftmargin=2em]
    \item an \emph{analytic} viewpoint, where one compares integrals of test
    functions against \(\mu\) and \(\nu\);
    \item a \emph{probabilistic} viewpoint, where one asks whether there are
    random variables \(X\sim\mu\) and \(Y\sim\nu\) on the same probability space
    satisfying a pathwise constraint, such as \(X\leq Y\) almost surely.
\end{itemize}

Optimal transport enters because a joint law of \((X,Y)\) with marginals
\(\mu\) and \(\nu\) is exactly a transport plan between \(\mu\) and \(\nu\).
Thus Strassen's theorem can be read as a feasibility theorem for constrained
transport plans.

\section{Couplings and transport plans}

Let \(E\) be a measurable space, and let \(\mu,\nu\in\Pcal(E)\) be probability
measures. A \emph{coupling} of \(\mu\) and \(\nu\) is a probability measure
\(\pi\in\Pcal(E\times E)\) such that
\[
        \pi(A\times E)=\mu(A),
        \qquad
        \pi(E\times A)=\nu(A)
\]
for every measurable set \(A\subseteq E\). The set of all couplings is denoted
by
\[
        \PiC(\mu,\nu).
\]
Equivalently, \(\pi\in\PiC(\mu,\nu)\) is the joint law of a pair of random
variables \((X,Y)\) such that
\[
        X\sim\mu,
        \qquad
        Y\sim\nu.
\]

In optimal transport, the same object is called a \emph{transport plan}. Given a
cost function \(c:E\times E\to (-\infty,\infty]\), the Kantorovich optimal
transport problem is
\[
        \mathcal T_c(\mu,\nu)
        :=
        \inf_{\pi\in\PiC(\mu,\nu)}
        \int_{E\times E} c(x,y)\,\pi(\dd x,\dd y).
\]
Strassen's theorem concerns the existence of couplings satisfying structural
constraints. This is a transport problem with constraints rather than only a
cost minimisation problem.

\section{Stochastic orders generated by test functions}

\begin{definition}[Order generated by test functions]
Let \(\mathcal F\) be a class of measurable real-valued functions on \(E\). We
write
\[
        \mu \preceq_{\mathcal F} \nu
\]
if
\[
        \int_E f\,\dd\mu
        \leq
        \int_E f\,\dd\nu
        \qquad
        \text{for every } f\in\mathcal F
\]
whenever the integrals are well-defined.
\end{definition}

Different choices of \(\mathcal F\) give different stochastic orders.

\subsection{Usual stochastic order}

Suppose \(E\) is partially ordered, with order relation denoted by \(\leq\). A
function \(f:E\to\R\) is \emph{increasing} if
\[
        x\leq y
        \quad\Longrightarrow\quad
        f(x)\leq f(y).
\]

\begin{definition}[Usual stochastic order]
For probability measures \(\mu,\nu\in\Pcal(E)\), one says that \(\mu\) is smaller
than \(\nu\) in the usual stochastic order, and writes
\[
        \mu\preceq_{\mathrm{st}}\nu,
\]
if
\[
        \int_E f\,\dd\mu
        \leq
        \int_E f\,\dd\nu
        \qquad
        \text{for every bounded increasing } f:E\to\R.
\]
\end{definition}

Intuitively, \(\mu\preceq_{\mathrm{st}}\nu\) means that a random variable with
law \(\nu\) lies to the right of, or above, a random variable with law \(\mu\).

\begin{example}[The real line]
On \(E=\R\) with the usual order,
\[
        \mu\preceq_{\mathrm{st}}\nu
\]
if and only if their distribution functions satisfy
\[
        F_\mu(t)\geq F_\nu(t)
        \qquad
        \text{for every } t\in\R.
\]
Equivalently, in terms of upper tails,
\[
        \mu([t,\infty))\leq \nu([t,\infty))
        \qquad
        \text{for every } t\in\R.
\]
The measure \(\nu\) has more mass to the right than \(\mu\).
\end{example}

\subsection{Convex order}

Another important order compares variability rather than location.

\begin{definition}[Convex order]
Let \(\mu,\nu\in\Pcal_1(\R^d)\), meaning that both measures have finite first
moment. One says that \(\mu\) is smaller than \(\nu\) in convex order, and writes
\[
        \mu\preceq_{\mathrm{cx}}\nu,
\]
if
\[
        \int_{\R^d} \varphi\,\dd\mu
        \leq
        \int_{\R^d} \varphi\,\dd\nu
\]
for every convex function \(\varphi:\R^d\to\R\) for which the two integrals are
well-defined.
\end{definition}

Since affine functions and their negatives are convex, convex order implies
that \(\mu\) and \(\nu\) have the same barycentre:
\[
        \int x\,\mu(\dd x)
        =
        \int y\,\nu(\dd y).
\]
The interpretation is that \(\nu\) is a mean-preserving spread of \(\mu\).

\begin{example}
On \(\R\),
\[
        \delta_0
        \preceq_{\mathrm{cx}}
        \frac12\delta_{-1}+\frac12\delta_1.
\]
Indeed, Jensen's inequality gives
\[
        \varphi(0)
        \leq
        \frac12\varphi(-1)+\frac12\varphi(1)
\]
for every convex \(\varphi\). The second law is more variable, but the mean is
still zero.
\end{example}

\section{Strassen's theorem for the usual stochastic order}

Strassen's theorem says that the integral comparison defining the usual
stochastic order is equivalent to the existence of an ordered coupling.

\begin{theorem}[Strassen theorem: usual stochastic order]
Let \(E\) be a Polish space equipped with a closed partial order \(\leq\). In
particular, one may take \(E=\R^d\) with the coordinatewise order. For
\(\mu,\nu\in\Pcal(E)\), the following are equivalent:
\begin{enumerate}[label=\textup{(\roman*)}, leftmargin=2em]
    \item \(\mu\preceq_{\mathrm{st}}\nu\), i.e.
    \[
            \int f\,\dd\mu
            \leq
            \int f\,\dd\nu
    \]
    for every bounded increasing function \(f\).
    \item There exists a coupling \(\pi\in\PiC(\mu,\nu)\) such that
    \[
            \pi\bigl(\{(x,y):x\leq y\}\bigr)=1.
    \]
    Equivalently, there exist random variables \(X\sim\mu\) and \(Y\sim\nu\)
    such that
    \[
            X\leq Y
            \qquad\text{almost surely.}
    \]
\end{enumerate}
\end{theorem}

The implication \((\mathrm{ii})\Rightarrow(\mathrm{i})\) is immediate. If
\(X\leq Y\) almost surely and \(f\) is increasing, then
\[
        f(X)\leq f(Y)
        \qquad\text{almost surely,}
\]
and therefore
\[
        \int f\,\dd\mu
        =
        \E[f(X)]
        \leq
        \E[f(Y)]
        =
        \int f\,\dd\nu.
\]

The nontrivial part is \((\mathrm{i})\Rightarrow(\mathrm{ii})\). It says that
all scalar inequalities against increasing test functions can be realised by a
single joint construction in which the lower random variable is below the upper
one almost surely.

\begin{example}[Quantile coupling on \(\R\)]
Let \(F_\mu^{-1}\) and \(F_\nu^{-1}\) be the quantile functions
\[
        F_\mu^{-1}(u)
        :=
        \inf\{x\in\R:F_\mu(x)\geq u\},
        \qquad 0<u<1.
\]
If \(\mu\preceq_{\mathrm{st}}\nu\), then
\[
        F_\mu^{-1}(u)
        \leq
        F_\nu^{-1}(u)
        \qquad
        \text{for every } u\in(0,1).
\]
Taking \(U\sim\mathrm{Unif}(0,1)\) and setting
\[
        X=F_\mu^{-1}(U),
        \qquad
        Y=F_\nu^{-1}(U),
\]
we get
\[
        X\sim\mu,
        \qquad
        Y\sim\nu,
        \qquad
        X\leq Y \quad\text{almost surely.}
\]
This is the monotone, or quantile, coupling.
\end{example}

\section{Strassen's theorem for convex order}

There is a second fundamental form of Strassen's theorem, where the pathwise
constraint \(X\leq Y\) is replaced by a conditional expectation constraint.

\begin{definition}[Martingale coupling]
Let \(\mu,\nu\in\Pcal_1(\R^d)\). A coupling \(\pi\in\PiC(\mu,\nu)\) is a
martingale coupling if, after disintegrating
\[
        \pi(\dd x,\dd y)=\mu(\dd x)\,\pi_x(\dd y),
\]
one has
\[
        \int_{\R^d} y\,\pi_x(\dd y)=x
        \qquad
        \text{for } \mu\text{-a.e. }x.
\]
Equivalently, if \((X,Y)\sim\pi\), then
\[
        \E[Y\mid X]=X.
\]
\end{definition}

\begin{theorem}[Strassen theorem: convex order]
Let \(\mu,\nu\in\Pcal_1(\R^d)\). Then
\[
        \mu\preceq_{\mathrm{cx}}\nu
\]
if and only if there exists a martingale coupling
\[
        \pi\in\PiC(\mu,\nu)
\]
between \(\mu\) and \(\nu\). Equivalently, there are random variables
\(X\sim\mu\) and \(Y\sim\nu\) such that
\[
        \E[Y\mid X]=X.
\]
\end{theorem}

Again, one implication follows from Jensen's inequality. If \(\E[Y\mid X]=X\)
and \(\varphi\) is convex, then
\[
        \varphi(X)
        =
        \varphi\bigl(\E[Y\mid X]\bigr)
        \leq
        \E[\varphi(Y)\mid X].
\]
Taking expectations gives
\[
        \int \varphi\,\dd\mu
        \leq
        \int \varphi\,\dd\nu.
\]
The converse is the martingale form of Strassen's theorem: the family of convex
integral inequalities is exactly the condition needed for the existence of a
martingale transport plan.

\section{Connection with optimal transport}

The link with optimal transport is simple but important: a coupling is a
transport plan. Stochastic orders ask whether there exists a transport plan with
special structure.

\subsection{Order-constrained transport}

Let
\[
        R:=\{(x,y)\in E\times E:x\leq y\}
\]
be the order relation. Define
\[
        \PiC_R(\mu,\nu)
        :=
        \{\pi\in\PiC(\mu,\nu):\pi(R)=1\}.
\]
Then Strassen's theorem says
\[
        \PiC_R(\mu,\nu)\neq\varnothing
        \qquad\Longleftrightarrow\qquad
        \mu\preceq_{\mathrm{st}}\nu.
\]
Thus the usual stochastic order is precisely the feasibility condition for a
transport plan supported on the order relation.

This can be expressed as an optimal transport problem by using the cost
\[
        c_R(x,y)
        :=
        \begin{cases}
        0, & x\leq y,\\
        +\infty, & x\not\leq y.
        \end{cases}
\]
Then
\[
        \inf_{\pi\in\PiC(\mu,\nu)}
        \int c_R(x,y)\,\pi(\dd x,\dd y)
\]
is finite if and only if there exists a coupling supported on \(R\), which is
exactly the conclusion of Strassen's theorem.

More generally, once feasibility is known, one may optimise a cost among only
ordered couplings:
\[
        \inf_{\pi\in\PiC_R(\mu,\nu)}
        \int c(x,y)\,\pi(\dd x,\dd y).
\]
This is an order-constrained optimal transport problem.

\subsection{Kantorovich duality viewpoint}

The transport formulation also explains why test functions appear. For a cost
\(c\), the Kantorovich dual problem has the form
\[
        \sup_{\phi,\psi}
        \left\{
        \int \phi\,\dd\mu+
        \int \psi\,\dd\nu
        :
        \phi(x)+\psi(y)\leq c(x,y)
        \right\}.
\]
For the constraint cost \(c_R\), the inequality
\[
        \phi(x)+\psi(y)\leq c_R(x,y)
\]
only imposes a restriction when \(x\leq y\). Taking
\[
        \phi=f,
        \qquad
        \psi=-f,
\]
the dual constraint becomes
\[
        f(x)-f(y)\leq 0
        \qquad
        \text{whenever } x\leq y,
\]
which means precisely that \(f\) is increasing. The corresponding dual value is
\[
        \int f\,\dd\mu-
        \int f\,\dd\nu.
\]
Thus increasing test functions are the dual certificates for the existence, or
non-existence, of an ordered coupling. This is the optimal-transport shadow of
Strassen's theorem.

\subsection{Martingale optimal transport}

For convex order, the constrained set of transport plans is
\[
        \PiC_M(\mu,\nu)
        :=
        \left\{
        \pi\in\PiC(\mu,\nu):
        \int y\,\pi_x(\dd y)=x\text{ for }\mu\text{-a.e. }x
        \right\}.
\]
Strassen's theorem says
\[
        \PiC_M(\mu,\nu)\neq\varnothing
        \qquad\Longleftrightarrow\qquad
        \mu\preceq_{\mathrm{cx}}\nu.
\]
Therefore convex order is the feasibility condition for \emph{martingale
optimal transport}. Given a cost \(c\), the martingale optimal transport problem
is
\[
        \inf_{\pi\in\PiC_M(\mu,\nu)}
        \int c(x,y)\,\pi(\dd x,\dd y),
\]
or similarly the corresponding supremum problem.

The martingale constraint changes the dual problem. A typical superhedging dual
for the supremum problem is
\[
        \sup_{\pi\in\PiC_M(\mu,\nu)}
        \int c(x,y)\,\pi(\dd x,\dd y)
        =
        \inf_{\varphi,\psi,h}
        \left\{
        \int \varphi\,\dd\mu+
        \int \psi\,\dd\nu
        \right\},
\]
where the infimum is over functions satisfying
\[
        \varphi(x)+\psi(y)+h(x)\cdot (y-x)
        \geq
        c(x,y)
        \qquad
        \text{for all }x,y.
\]
The additional term \(h(x)\cdot(y-x)\) is invisible under martingale couplings,
because
\[
        \int h(x)\cdot(y-x)\,\pi(\dd x,\dd y)
        =0
\]
whenever \(\E[Y\mid X]=X\). This is the same mechanism by which affine
functions do not distinguish two measures in convex order once their means are
equal.

\section{Comparison of the main orders}

\begin{center}
\begin{tabular}{c|c|c|c}
Order & Test functions & Coupling condition & OT interpretation \\
\hline
Usual stochastic order & increasing \(f\) & \(X\leq Y\) a.s. & support-constrained OT \\
Convex order & convex \(\varphi\) & \(\E[Y\mid X]=X\) & martingale OT \\
Increasing convex order & increasing convex \(\varphi\) & submartingale type constraint & submartingale OT
\end{tabular}
\end{center}

The common pattern is the following:
\[
        \text{integral inequalities against test functions}
        \quad\Longleftrightarrow\quad
        \text{existence of a constrained coupling}.
\]
Optimal transport provides the language and tools for studying such couplings.
Classical OT optimises over all couplings; Strassen-type theorems identify when
the admissible class remains nonempty after imposing order, martingale, or
similar structural constraints.

\section{Finite-state intuition}

On a finite partially ordered set \(E\), the usual stochastic order has a simple
network-flow interpretation. A coupling \(\pi\) is a nonnegative matrix with row
sums \(\mu\) and column sums \(\nu\). The ordered-coupling constraint says
\[
        \pi(x,y)=0
        \qquad
        \text{unless }x\leq y.
\]
Thus Strassen's theorem becomes a feasible-flow theorem: mass initially located
at \(x\) may only be transported upward to states \(y\geq x\). The condition
\(\mu\preceq_{\mathrm{st}}\nu\) says exactly that every increasing test function
has no larger expectation under \(\mu\) than under \(\nu\), or equivalently that
the target law \(\nu\) has enough mass in every upper region to receive the mass
that must be moved upward.

\section{Summary}

Stochastic orders compare probability measures by testing them against a class
of functions. Strassen's theorem converts these analytic comparisons into the
existence of couplings with structural constraints:
\[
        \mu\preceq_{\mathrm{st}}\nu
        \quad\Longleftrightarrow\quad
        \exists (X,Y): X\sim\mu,\ Y\sim\nu,\ X\leq Y\text{ a.s.},
\]
and
\[
        \mu\preceq_{\mathrm{cx}}\nu
        \quad\Longleftrightarrow\quad
        \exists (X,Y): X\sim\mu,\ Y\sim\nu,\ \E[Y\mid X]=X.
\]
Since a coupling is the same thing as a transport plan, these results are
naturally part of optimal transport. The usual stochastic order corresponds to
transport supported on an order relation, while convex order corresponds to
martingale transport. Kantorovich duality explains why the relevant test
functions appear as dual certificates for the feasibility of these constrained
transport problems.

\begin{thebibliography}{9}

\bibitem{Strassen1965}
V. Strassen,
\emph{The existence of probability measures with given marginals},
Annals of Mathematical Statistics, 36 (1965), 423--439.

\bibitem{ShakedShanthikumar}
M. Shaked and J. G. Shanthikumar,
\emph{Stochastic Orders},
Springer, 2007.

\bibitem{MuellerStoyan}
A. M\"uller and D. Stoyan,
\emph{Comparison Methods for Stochastic Models and Risks},
Wiley, 2002.

\bibitem{Villani}
C. Villani,
\emph{Optimal Transport: Old and New},
Springer, 2009.

\bibitem{BeiglbockJuillet}
M. Beiglb\"ock and N. Juillet,
\emph{On a problem of optimal transport under marginal martingale constraints},
Annals of Probability, 44 (2016), 42--106.

\end{thebibliography}

\end{document}
